% GENERATED FILE. Edit canonical lessons, then run npm run generate:books.
% canonical-source-hash: fnv1a64:cdf392821d8af2f4
% canonical-lessons: UR-C10-kaan, UR-C10-naak, UR-C10-aankh, UR-C10-munh

\chapter{Four Words on the Face}
\label{ch:ur-face-words}
\hlchaptermodality{10}

\begin{chapteropening}
\textbf{By the end of this chapter:} \emph{I can name my ear, nose, eye and mouth in Urdu, notice that kān and nāk share the same three letters reversed, and read the nūn that marks a nasal vowel instead of a spoken consonant.}

Name kān, nāk, āṅkh and muṅh, explain the plain nūn that nasalizes āṅkh and muṅh against the dedicated \ur{ں} of maiṅ and hūṅ, and count three settled Indo-European cousins found this chapter — nāk/nose, āṅkh/eye — against kān's honest dead end.
\end{chapteropening}

\section[kān]{\ur{کان} — \textquotedblleft{}ear,\textquotedblright{} where the body words start}

\label{lesson:UR-C10-kaan}



\begin{quote}
Say \textbf{merā khāndān} — my family, closing the last chapter. This chapter turns from the people around you to the wellbeing check itself: four words for the body, starting with the one you use to hear the question.
\end{quote}

\subsection*{You'll want to know first — one word}
\begin{quote}
\textbf{\ur{کان}} — \emph{kān} — \textbf{ear}
\end{quote}

\begin{quote}
\textbf{\ur{میرا} \ur{کان۔}} — \emph{merā kān.} — \textquotedblleft{}my ear.\textquotedblright{}
\end{quote}

\textbf{\ur{کان}} is masculine, so \textbf{merā} keeps the form you already know from \textbf{\ur{بھائی}} and \textbf{\ur{خاندان}}.

\begin{cousinweb}
\textbf{kān} continues Sanskrit \emph{karṇa} by way of Prakrit and Apabhramsa \emph{kaṇṇa}. That much is solid. Where \emph{karṇa} itself comes from is not: the trail runs out inside Indo-Aryan, with only a loose, unconfirmed suggestion linking it to Persian \textbf{\ur{کر}} (\emph{kar}), \textquotedblleft{}deaf.\textquotedblright{} No English cousin, no settled Indo-European root — the same honest dead end \textbf{\ur{پڑھنا}}'s \emph{paṭh-} hit two chapters ago.

Notice the pattern from last chapter still holds: \textbf{\ur{کان}} is inherited, straight from Sanskrit, the same road as \textbf{\ur{بھائی}} and \textbf{\ur{بہن}}. \textbf{\ur{خاندان}}, the word that closed that chapter, was Persian instead — the body's own words stay native even where their history goes quiet.
\end{cousinweb}

\subsection*{Your turn}
\begin{itemize}
\raggedright
  \item \emph{Say it:} \textbf{kān} — ear; then \textbf{merā kān}
  \item \emph{Say it:} \textbf{khāndān}, then \textbf{kān} — one word ends, the body words begin
  \item \emph{Name:} what's uncertain about \emph{karṇa}'s own origin
\end{itemize}

\begin{tcolorbox}[breakable,colback=teal!4,colframe=teal!35!black,title={Before you move on}]
Is \textbf{\ur{کان}} masculine or feminine? (\textbf{Masculine.}) Does \emph{karṇa} have a secure English cousin? (\textbf{No — the origin is unsettled.})

Sources: \href{https://en.wiktionary.org/wiki/\%DA\%A9\%D8\%A7\%D9\%86}{Wiktionary: \ur{کان}}.
\end{tcolorbox}
\section[nāk]{\ur{ناک} — \textquotedblleft{}nose,\textquotedblright{} the same three letters as \ur{کان}, backwards}

\label{lesson:UR-C10-naak}



\begin{quote}
Say \textbf{merā kān} — my ear, masculine, unchanged \textbf{merā}. This word shares every letter with it.
\end{quote}

\subsection*{You'll want to know first — one word}
\begin{quote}
\textbf{\ur{ناک}} — \emph{nāk} — \textbf{nose}
\end{quote}

\begin{quote}
\textbf{\ur{میری} \ur{ناک۔}} — \emph{merī nāk.} — \textquotedblleft{}my nose.\textquotedblright{}
\end{quote}

\textbf{\ur{ناک}} is feminine — \textbf{merī}, not \textbf{merā}. Grammatical gender does not follow shape or sound; it has to be learned noun by noun.

\subsection*{The letters in this word}
\textbf{\ur{کان}} is \textbf{\ur{ک} \ur{ا} \ur{ن}} — \emph{k, ā, n}. \textbf{\ur{ناک}} is \textbf{\ur{ن} \ur{ا} \ur{ک}} — \emph{n, ā, k}. The same three letters, in exactly the opposite order. Ear and nose, spelled as mirror images of each other — a coincidence of Urdu's script, not of its history; the two words are not related.

\begin{cousinweb}
\textbf{nāk} continues Sauraseni Prakrit \emph{ṇakka}, from Proto-Indo-Iranian \textbf{*náHs}, from Proto-Indo-European \textbf{*néh\textsubscript{2}s}, \textquotedblleft{}nose.\textquotedblright{} English \textbf{nose} comes from that very root. Where \textbf{\ur{کان}}'s ancestry ran out unanswered — recall \emph{karṇa}'s own origin was left unsettled last lesson — \textbf{\ur{ناک}}'s does not: \emph{nāk} and \emph{nose} are genuine cousins, joining \textbf{\ur{بھائی}} and \emph{brother} as this book's settled Indo-European matches.
\end{cousinweb}

\subsection*{Your turn}
\begin{itemize}
\raggedright
  \item \emph{Say it:} \textbf{nāk} — nose; then \textbf{merī nāk}
  \item \emph{Spell:} \textbf{kān} forward, then \textbf{nāk} — the same three letters, reversed
  \item \emph{Connect:} \textbf{nāk} $\leftarrow$ \emph{*néh\textsubscript{2}s} $\to$ English \textbf{nose}, real cousins
  \item \emph{Contrast:} \textbf{merā kān} — \textbf{merī nāk}
\end{itemize}

\begin{tcolorbox}[breakable,colback=teal!4,colframe=teal!35!black,title={Before you move on}]
What is the letter relationship between \textbf{\ur{کان}} and \textbf{\ur{ناک}}? (\textbf{Same three letters, reversed order.}) Is \textbf{\ur{ناک}}'s English cousin settled or disputed? (\textbf{Settled — \emph{nose} shares its root.}) Is \textbf{\ur{ناک}} masculine or feminine? (\textbf{Feminine.})

Sources: \href{https://en.wiktionary.org/wiki/\%D9\%86\%D8\%A7\%DA\%A9}{Wiktionary: \ur{ناک}}.
\end{tcolorbox}
\section[āṅkh]{\ur{آنکھ} — \textquotedblleft{}eye,\textquotedblright{} and a fifth letter learns to breathe}

\label{lesson:UR-C10-aankh}



\begin{quote}
Say \textbf{merī nāk} — my nose, feminine. The next word on the face is feminine too, and it puts the two-eyed \textbf{\ur{ھ}} on a letter it has never touched.
\end{quote}

\subsection*{You'll want to know first — one word}
\begin{quote}
\textbf{\ur{آنکھ}} — \emph{āṅkh} — \textbf{eye}
\end{quote}

\begin{quote}
\textbf{\ur{میری} \ur{آنکھ۔}} — \emph{merī āṅkh.} — \textquotedblleft{}my eye.\textquotedblright{}
\end{quote}

\subsection*{The letters in this word}
From the right edge: \textbf{\ur{آ}}, alif carrying long \emph{ā} (the same mark you already read in \textbf{\ur{آپ}}), then \textbf{\ur{ن}}, then \textbf{\ur{ک}} \emph{k}, then \textbf{\ur{ھ}}.

\textbf{\ur{کھ}} extends the two-eyed \emph{he} to a fifth base letter. It has already aspirated \textbf{\ur{ج}} in \emph{samajh-}, \textbf{\ur{ٹ}} in \emph{ṭhīk}, \textbf{\ur{ڑ}} in \emph{paṛh-}, and \textbf{\ur{چ}} in \emph{pūchh-}; here it does the same job to \textbf{\ur{ک}}, giving \emph{kh}.

The \textbf{\ur{ن}} in this word is not pronounced as its own consonant. It marks the \textbf{\ur{آ}} before it as nasal — the vowel said through the nose, with no separate \emph{n} sound following it. Say \textbf{āṅkh} by nasalizing the \emph{ā}, not by adding a clean \emph{n}.

\begin{cousinweb}
\textbf{āṅkh} continues Sanskrit \emph{ákṣi}, from Proto-Indo-European \textbf{*h\textsubscript{3}ek\textsuperscript{w}-}, \textquotedblleft{}eye, to see.\textquotedblright{} English \textbf{eye} descends from that same root. This chapter has now found two settled Indo-European cousins in a row — \textbf{\ur{ناک}}/\emph{nose}, \textbf{\ur{آنکھ}}/\emph{eye} — after \textbf{\ur{کان}}'s honest dead end.
\end{cousinweb}

\subsection*{Your turn}
\begin{itemize}
\raggedright
  \item \emph{Say it:} \textbf{āṅkh} — eye; then \textbf{merī āṅkh}
  \item \emph{Nasalize:} \textbf{ā} through the nose, no separate \textbf{n}
  \item \emph{Count:} five letters now carrying \textbf{\ur{ھ}} — \textbf{\ur{ج،} \ur{ٹ،} \ur{ڑ،} \ur{چ،} \ur{ک}}
  \item \emph{Connect:} \textbf{āṅkh} $\leftarrow$ \emph{*h\textsubscript{3}ek\textsuperscript{w}-} $\to$ English \textbf{eye}
\end{itemize}

\begin{tcolorbox}[breakable,colback=teal!4,colframe=teal!35!black,title={Before you move on}]
Which letter does \textbf{\ur{کھ}} newly pair the two-eyed \emph{he} with, and name two letters it already paired with? (\textbf{\ur{ک}}; any two of \textbf{\ur{ج،} \ur{ٹ،} \ur{ڑ،} \ur{چ}}.) What job does \textbf{\ur{ن}} do in \textbf{\ur{آنکھ}}? (\textbf{Marks the \ur{آ} as nasal — it is not its own consonant.}) Is \textbf{\ur{آنکھ}}'s English cousin settled? (\textbf{Yes — \emph{eye}.})

Sources: \href{https://en.wiktionary.org/wiki/\%D8\%A2\%D9\%86\%DA\%A9\%DA\%BE}{Wiktionary: \ur{آنکھ}}.
\end{tcolorbox}
\section[muṅh]{\ur{منہ} — \textquotedblleft{}mouth,\textquotedblright{} and the \ur{نون} \ur{آنکھ} already used without explaining}

\label{lesson:UR-C10-munh}



\begin{quote}
Three face words: \textbf{kān, nāk, āṅkh} — recall that \textbf{\ur{کان}} and \textbf{\ur{ناک}} shared the same three letters, reversed. One more word closes the chapter, and it finally explains something \textbf{\ur{آنکھ}} already did.
\end{quote}

\subsection*{You'll want to know first — one word}
\begin{quote}
\textbf{\ur{منہ}} — \emph{muṅh} — \textbf{mouth} (also: \textbf{face})
\end{quote}

\begin{quote}
\textbf{\ur{میرا} \ur{منہ۔}} — \emph{merā muṅh.} — \textquotedblleft{}my mouth.\textquotedblright{}
\end{quote}

\textbf{\ur{منہ}} is masculine.

\subsection*{The letters in this word}
From the right edge: \textbf{\ur{م}} \emph{m}, then \textbf{\ur{ن}}, then \textbf{\ur{ہ}}, the round \emph{he} you already write in \textbf{\ur{ہوں}} and \textbf{\ur{ہونا}} — a real, spoken \emph{h}, not the two-eyed aspirator.

That middle \textbf{\ur{ن}} does the same job it did in \textbf{\ur{آنکھ}}, unglossed at the time: it marks the vowel before it as nasal, not a separate \emph{n} sound. Say \textbf{muṅh} by nasalizing the \emph{u}, with no clean \emph{n} in between.

Keep this apart from \textbf{\ur{میں}} \emph{maiṅ} and \textbf{\ur{ہوں}} \emph{hūṅ}, whose nasal ending is written with a different letter entirely: \textbf{\ur{ں}}, \emph{nūn ghunna}, a dedicated nasalization mark. \textbf{\ur{منہ}} and \textbf{\ur{آنکھ}} nasalize with a plain \textbf{\ur{ن}} doing double duty; \textbf{\ur{میں}} and \textbf{\ur{ہوں}} nasalize with \textbf{\ur{ں}}, a letter built for exactly that job and nothing else. Same sound, two different letters, depending on which word you are in.

\begin{cousinweb}
\textbf{muṅh} continues Sanskrit \emph{múkha}, \textquotedblleft{}mouth, face.\textquotedblright{} One proposal, from the Sanskrit etymologist Manfred Mayrhofer, ties it to a Proto-Indo-European gesture-root \textbf{*mu-}, \textquotedblleft{}pressed lips\textquotedblright{} — the same root that would sit behind Greek \emph{mŷthos} (\textquotedblleft{}word,\textquotedblright{} behind English \textbf{myth}) and German \emph{Maul}, \textquotedblleft{}mouth.\textquotedblright{} Mayrhofer himself calls a non-Indo-European origin less likely but does not rule it out. Hold this one the way you held \textbf{\ur{بہن}}'s link to \emph{bhaj-} — a real proposal, not a verdict.
\end{cousinweb}

\subsection*{Your turn}
\begin{itemize}
\raggedright
  \item \emph{Say it:} \textbf{muṅh} — mouth; then \textbf{merā muṅh}
  \item \emph{Contrast:} plain \textbf{\ur{ن}} nasalizing in \textbf{āṅkh} and \textbf{muṅh}; \textbf{\ur{ں}} \emph{nūn ghunna} nasalizing in \textbf{maiṅ} and \textbf{hūṅ}
  \item \emph{Say it:} four face words now — \textbf{kān, nāk, āṅkh, muṅh}
  \item \emph{Connect:} \textbf{muṅh} possibly $\leftarrow$ \emph{*mu-} $\to$ Greek \textbf{mŷthos} $\to$ English \textbf{myth}, held loosely
\end{itemize}

\begin{tcolorbox}[breakable,colback=teal!4,colframe=teal!35!black,title={Before you move on}]
Which earlier word in this chapter used a plain \textbf{\ur{نون}} for nasalization the same way \textbf{\ur{منہ}} does? (\textbf{\ur{آنکھ}}.) Which letter does \textbf{\ur{میں}} use instead? (\textbf{\ur{ں}}, \emph{nūn ghunna}.) Is \textbf{\ur{منہ}}'s link to Greek \emph{mŷthos} certain? (\textbf{No — a proposal, held loosely.})

Sources: \href{https://en.wiktionary.org/wiki/\%D9\%85\%D9\%86\%DB\%81}{Wiktionary: \ur{منہ}}.
\end{tcolorbox}
